\ifdefined\ThesisLanguageIsEnglish
\chapter*{Abstract}
\else
\chapter*{Kurzfassung}
\fi
\label{ch:Abstract}

Diese Arbeit entwickelt ein unüberwachtes Verfahren zur Erkennung auffälliger
Passanten- und Verweildauerdaten aus Bad Nauheim. Die Datenaufbereitung fasst
Mehrfachmessungen zusammen, kennzeichnet Ausfälle und ergänzt interne Lücken,
ohne Informationen aus dem späteren Testzeitraum zu verwenden. Ein einfaches
LSTM sagt aus den vorhergehenden 24 Stunden die nächste Stunde voraus. Große
Abweichungen zwischen Vorhersage und Messung gelten als mögliche Anomalien.
Unter gleichen Testbedingungen benötigt das aktuelle LSTM 5.170 statt 17.282
Parameter. Bei starken künstlichen Anomalien steigt der mittlere F1-Wert
gegenüber dem Basismodell von 0{,}401 auf 0{,}491. Die normale
Besucherprognose wird dabei etwas schlechter. Der Vorteil liegt daher in der
Erkennung der geprüften Anomalien und nicht in jeder einzelnen Kennzahl. Die
nur zweiwöchige Datenbasis und fehlende vollständige reale Markierungen
begrenzen die Aussagekraft.
